\section{Rapport : 17/06/2026}

Durant la réunion on a proposé d'affiner les hypothèses naturelles que l'on pouvait supposer
pour trouver une solution à l'équation~\ref{eq:ev_nonop}. Enfin, on a également réussi à démontrer
que l'unique solution $u \in C^1([0,T], L^1)$ vérifiait la proprité importante
\[
	\forall t \in [0,T], \quad u(t) \ge 0,\quad \lambda\ p.p.
\]
Pour démontrer ces résultats, il est intéressant de rappeler des faits générals sur les problèmes de
Cauchy.

\subsection{Problème de Cauchy}

Soit $E$, un espace de Banach et $f : U \to E$ tel que $U \subset \R \times E$ est un ouvert. Un problème de Cauchy
se décrit par l'équation

\begin{equation}
	\begin{cases}
		v'(t) = f(t,v(t)) \\
		v(t_0) = u_0,
	\end{cases}
	\label{eq:cauchy}
\end{equation}
avec $u_0 \in E$ et $t_0 \in \R$. On verra alors que si $f$ vérifie certaines conditions de régularité,
alors le problème admet une unique solution maximale \ref{defn:cauchy}.

\begin{defn}{Problème de Cauchy}{cauchy}
	\begin{itemize}
		\item On dit que $(u,I)$ est une solution au problème de Cauchy \ref{eq:cauchy} si
		      \[
			      \begin{cases}
				      \forall t \in I, \ u'(t) = f(t,u(t)) \\
				      \forall t \in I, \ (t,u(t)) \in U    \\
				      t_0 \in I, \quad u(t_0) = u_0.
			      \end{cases}
		      \]
		      avec $I$ un intervalle ouvert de $\R$.
		\item On dit que deux solutions $(u, I) \le (v,J)$ si $I \subset J$ et si $v_{|I} = u$. C'est une relation d'ordre sur l'ensemble
		      des solutions au problème de Cauchy.
		\item Une solution maximale $(u,I)$ est une solution qui est maximale pour la relation d'ordre définie précédemment.
	\end{itemize}
\end{defn}

L'existence et l'unicité de solutions dépendent principalement de la régularité de la fonction $f$. Par exemple,
si $f$ est semi-lipschitzienne et continue, alors le problème de Cauchy admet une unique solution maximale.

\begin{defn}{semi-lipschitzianité}{}
	La fonction $f$ est semi-lipschitzienne si pour tout $(t_0, x_0)\in U$, il existe une constante $c > 0$, un voisinage ouvert $I$ qui contienne $t_0$ et
	un voisinage ouvert $V$ qui contienne $x_0$ tel que $I \times V \subset U$ et
	\[
		\forall t \in I,\forall x,y \in V, \quad \|f(t,x) - f(t,y)\| \le c \|x-y\|.
	\]
\end{defn}

Les résultats suivants viennent de \cite{paulin}.

\begin{thm}{}{}
	\begin{itemize}
		\item Si $f$ est semi-lipschitzienne et continue, alors le problème de Cauchy \ref{eq:cauchy} admet une unique solution maximale $(u_{t_0, u_0}, I_{t_0, u_0})$.
		\item si $f$ est de classe $C^k$ alors une solution $u$ est de classe $C^{k+1}$.
	\end{itemize}
\end{thm}

\begin{defn}{Résolvante}{}
	Si pour tout $(t, x) \in U$, le problème de Cauchy admet une unique solution maximale $(u_{t,x}, I_{t,x})$, on peut poser $\aD = \left\{(s,t,x) \ :\ s \in I_{t,x} \right\}$ et définir la résolvante
	\[
		\aR : \aD \to E, \quad (s,t,x) \mapsto u_{t,x}(s).
	\]
	Dans ce cas, on peut également définir les ensembles $U_{s,t} = \{x \in E \ : \ (s,t,x) \in \aD\}$ pour tous $s,t \in \R$ et l'application
	\[
		\aR_{s,t} : U_{s,t} \to E, \quad x \mapsto u_{t,x}(s)
	\]
\end{defn}

\begin{thm}{}{}
	Si $f$ est une application semi-lipschitzienne et continue, alors $\aD$ est ouvert et pour tout $(s_0, t_0, x_0) \in \aD$, il existe un voisinage ouvert
	$V \subset \aD$ de $(s_0, t_0, x_0)$ tel que $\aR$ est lipschitzienne sur $V$.
\end{thm}

La notion de résolvante est particulièrement intéressant pour le cas linéaire du problème de Cauchy
\begin{equation}
	\begin{cases}
		v'(t) = A(t) \cdot v(t) + b(t) \\
		v(0) = u_0
	\end{cases}
	\label{eq:cauchy_lin}
\end{equation}
Avec $A : I \to \aL(E)$ et $b : I \to E$ continues. Dans ce cas l'unique solution maximale est définie sur $I$ tout entier.
Pour résudre ce problème, on considère le cas plus simple où
\[
	v'(t) = A(t) \cdot v(t),
\]
et on note $\aR$ la résolvante de cette équation.

\begin{prop}{}{}
	\begin{enumerate}
		\item $\forall t,s \in I, \quad \aR_{s,t} \in \aG\aL(E)$
		\item $\forall \alpha, s, t \in I, \quad \aR_{s,\alpha} \circ \aR_{\alpha, t} = \aR_{s,t}$.
		\item $\forall t \in I$, on définit
		      \[
			      Y : I \to \aL(E), \quad Y(s) \mapsto \aR_{s,t}.
		      \]
		      $Y$ est la solution maximale de l'équation
		      \[
			      \frac{dY}{ds} = Y \circ A(s), \quad Y(0) = Id_E.
		      \]
	\end{enumerate}
\end{prop}

Cette interprétation est très pratique car elle permet de résoudre le problème de Cauchy affine (avec $b(t) \neq 0$) en utilisant la méthode
de variaion de la constante.

\begin{prop}{}{}
	Le problème de Cauchy affine admet une unique solution maximale $(u, I)$ telle que
	\[
		\forall t \in I, \quad u(t) = \aR_{t,t_0} \cdot u_0 + \int_{t_0}^t \aR_{t,s}\cdot b(s) ds.
	\]
\end{prop}

\subsection{Résolution de l'équation de Focker--Planck (cas non linéaire)}

On cherche à résoudre l'équation \ref{eq:ev_nonop}. Pour cela on définit
\[
	f(w,v)(x) = \int g(y,x) F(y,w) v(y) dy, \quad h(w,v)(x) = -F(x,w)v(x)
\]
et
\[
	\bar{f}(v) = f(v,v), \quad \bar{h}(v) = h(v,v).
\]
Le problème \ref{eq:ev_nonop}, se reformule donc
\begin{equation}
	\begin{cases}
		v'(t) = \bar{f}(v(t)) + \bar{h}(v(t)) \\
		v(0) = u_0
	\end{cases}
	\label{eq:ev_new}
\end{equation}

\begin{prop}{}{}
	On suppose que $F : L^1 \to L^\infty$ est localement lipschitzienne. Alors,
	\begin{enumerate}
		\item $f : L^1 \times L^1 \to L^1$ et $h : L^1 \times L^1 \to L^1$.
		\item $\bar f$ et $\bar h$ sont localement lipschitziennes.
		\item $\forall v \in L^1, \ $ $\aA(v) = h(v, \cdot) \in \aL(L^1)$.
		\item $v \mapsto \aA(v)$ est continue.
	\end{enumerate}
\end{prop}

Le fait que $\bar f$ et $\bar h$ soient localement lipschitziennes impliquent nécessairement l'existence d'une unique solution maximale
$(u,I)$ telle que $u \in C^1(I, L^1)$ pour le problème \ref{eq:ev_new}. Pour la suite on considérera la restriction de $I$ à $\R_+$. D'où, on considère que $I$ est de la forme $[0, \alpha[$.
De plus, si on pose $A(t) = \aA(u(t))$ et $b(t) = \bar f (u(t))$, alors
$(u,I)$ est l'unique solution maximale du problème
\begin{equation}
	\begin{cases}
		v'(t) = b(t) + A(t) \cdot v(t) \\
		v(0) = u_0.
	\end{cases}
	\label{eq:ev_lin}
\end{equation}
Donc, si on note la résolvante $\aR^A$ pour le problème
\[
	v'(t) = A(t) \cdot v(t),
\]
alors,
\[
	\forall t \in I, \quad u(t) = \aR^A_{t,0}(u_0) + \int_{0}^t \aR^A_{t,s} \cdot b(s) ds
	= \aR^A_{t,0}(u_0) + \int_{0}^t \aR^A_{t,s} \cdot \bar f (u(s)) ds
\]
On peut donc prouver le résultat suivant :

\begin{prop}{}{}
	Si $(u, I)$ est la solution maximale du problème \ref{eq:ev_nonop} où $u_0 \ge 0$, alors il existe $\varepsilon > 0$ tel $\forall t \in [0,\varepsilon[, \ u(t) \ge 0$.
\end{prop}

\begin{proof}
	Si $(u, I)$ est la solution maximale et que $[0,T] \subset I$, alors u est un point fixe de l'application
	\[
		\Phi_T^0 : C([0,T], L^1) \to C([0,T], L^1)
	\]
	définie par
	\[
		\Phi_T^0(v)(t) = \aR^A_{t,0}(u_0) + \int_{0}^t \aR^A_{t,s} \cdot \bar f (v(s)) ds, \quad \forall t \in [0,T].
	\]
	On rappelle que $\aR^A$ et $\bar f$ sont localement lipschitziennes. Par conséquent, si $T$ est suffisament proche de $0$ et si $v, w$ sont suffisament proches de $u_0$, alors
	\[
		\|\Phi_T^0(w) -
		\Phi_T^0(v)\| \le T K \|w-v\|.
	\]
	Donc, si $T$ est assez petit, alors l'application $\Phi_T^0$ est contractante. De plus, on peut également choisir $T$ correctement pour que
	\[
		\Phi_T^0 : C([0,T], \bar{B}(u_0, r)) \to C([0,T], \bar{B}(u_0, r))
	\]
	et
	\[
		\forall t \in [0,T], \quad u(t) \in  \bar{B}(u_0, r).
	\]
	Alors, l'application $u$ restreinte à $[0,T]$ est l'unique point fixe de $\Phi_T^0$.
	Enfin, on remarque que l'ensemble $L^1_+ = \{v \in L^1 \ : \ v \ge 0\}$ est un fermé de $L^1$ et on note $B_+ = L^1_+ \cap \bar{B}(u_0, r)$.
	Pour conclure, on suppose que $\Phi_T^0$ envoie des éléments $C([0,T], L^1_+)$ dans $C([0,T], L^1_+)$. Par conséquent, l'application
	\[
		\Phi_T^1 : C([0,T], B_+) \to C([0,T], B_+)
	\]
	(vue comme la restriction de $\Phi_T^0$) est contractante et admet un point fixe $v$.

	\medskip

	On conclut en remarquant que $v$ est également un point fixe pour $\Phi_T^0$.
	Or, l'application $\Phi_T^0$ admet pour unique point fixe l'application $u$. Donc les deux solutions $u$ et $v$ sont identiques. D'où pour tout $t \in [0,T], \ u(t) \in L^1_+$.
\end{proof}

\begin{lem}{}{}
	$L^1_+ = \{v \in L^1 \ : \ v \ge 0\}$ est fermé dans $L^1$.
\end{lem}

\begin{proof}
	Soit $f_n \in L^1_+$ tel que $f_n \to f \in L^1$ quand $n \to +\infty$. Alors, il existe une sous-suite $(f_{\varphi(n)})$ de $(f_n)$ telle que
	\[
		f_{\varphi(n)} \to f \quad \lambda \ p.p.
	\]
	On note $C$ l'ensemble des points où il y a convergence et $P_n$ l'ensemble des points $x$ où $f_n(x) \ge 0$. En posant
	$A = C \cap \left(\cap_{n \in \N}P_n \right)$,
	on obtient que
	\[
		\forall x \in A \quad f(x) \ge 0, \quad \lambda(A^c) = 0.
	\]
	Donc $f \in L^1_+$.
\end{proof}

\begin{lem}{}{pos}
	Soit $T > 0$ tel que $[0,T] \subset I$, alors
	\[
		\Phi_T : C([0,T], L^1_+) \to C([0,T], L^1_+).
	\]
\end{lem}

\begin{proof}
	Pour prouver ce résultat, on va directement déterminer la résolvante $\aR^A$. On rappelle que cette résolvante est associée au problème
	\begin{equation}
		v'(t) = - \alpha(t) v(t)
		\label{eq:ev_fin}
	\end{equation}
	avec $\alpha(t) = F(\cdot, u(t)) \in L^{\infty}$. On sait que
	\[
		\alpha : [0,T] \to L^\infty
	\]
	et que $\alpha$ est continue. Par conséquent, comme $L^\infty$ est une algèbre de Banach commutative, on peut définir
	\[
		\psi(t,s) = \exp\left(-\int_s^t \alpha(s) ds \right) \in L^\infty.
	\]
	Alors, l'application $\psi(\cdot , s) : [0,T] \to L^\infty$ est de classe $C^1$.
	On montre sans problème que la résolvante
	\[
		\aR^A_{s,t} = \psi(s,t).
	\]
	où $L^\infty$ peut être vu comme un sous-ensemble de $\aL(L^1)$.
	Pour conclure, on remarque que $R_{s,t}^A : L^1_+ \to L^1_+$, que $\bar f : L^1_+ \to L^1_+$ et que
	l'intégrale de Bochner préserve également la positivité. On en conclut que
	\[
		\Phi_T(v)(t) = \aR^A_{t,0}(u_0) + \int_{0}^t \aR^A_{t,s} \cdot \bar f (v(s)) ds \in L^1_+, \quad \forall t \in [0,T].
	\]
\end{proof}

Avec ces deux lemmes, on vient de prouver que la positivité est une propriété qui se préserve localement à droite.
Avec un raisonnement simple, on peut montrer que cette propriété se préserve sur tout l'intervalle $I$.

\begin{rem}
	Pour démontrer le lemme \ref{lem:pos}, on a fait l'hypothèse supplémentaire que si $u \in L^1_+$, alors $F(\cdot , u) : L^1_+ \to L^1_+$. Cela est cohérent avec la définition de l'article \cite{geo} qui donne une formule explicite pour le taux de fragmentation \ref{eq:dens}.
\end{rem}

\begin{prop}{}{}
	Soit $(u,I)$, l'unique solution au problème de Cauchy \ref{eq:ev_nonop}, alors pour tout $t \in I$, $\quad u(t) \in L^1_+$.
\end{prop}

\begin{proof}
	On pose $J = \{t \in I \ : \ u(t) \notin L^1_+\}$ un ouvert de $I$ et on le suppose non vide. On note $0 \le m = \inf J$. Comme $0 \notin J$, alors $m \notin J$.
	On en déduit que $u(m) \in L^1_+$. Donc d'après le résultat précédent, il existe un $\varepsilon > 0$ tel que pour tout $0 \le  t < \varepsilon$,
	$u(m + t) \in L^1_+$. Cela contredit que $m$ est une borne inférieur pour $J$. On en déduit que $J$ est vide et que pour tout $t\in I$, $u(t) \in L^1_+$.
\end{proof}

\subsection{Introduction aux processus de Markov non-homogènes}

On a également pu présenter des éléments concernant les processus de Markov généraux et en particulier, non-homogènes. En effet, dans l'article \cite{frag_deaconu}, l'équation
de mesure introduisait naturellement un générateur infinitésimal pour la génération d'une chaîne de Markov. Cela ne peut pas fonctionner ici car le potentiel candidat pour le \og générateur infinitésimal \fg{} n'est pas linéaire.
Cela rend donc impossible toute construction d'un processus de Markov homogène caractérisé par un semi-groupe de transition.
Néanmoins, cela n'empêche pas a priori de construire un processus de Markov qui décrive l'équation de mesure. D'après le livre de Blumenthal \cite{blumenthal}, on a donc rappelé la définition générale d'un processus de Markov.

\begin{defn}{}{}
	Soit $(\Omega, \aF, (\aF_t)_{t \ge 0}, \mathbb{P})$ un espace probabilisé filtré et $(X_t)_{t \ge 0}$ un processus adapté à valeurs dans $(E, \aE)$. On dit que le processus $X$ est markovien
	si pour tout $t \ge 0$, les tribus $\aF_t$ et $\aF^t =\sigma(X_s : s \ge t)$ sont indépendantes conditionnellement à $X_t$. Cela se traduit par :
	\[
		\forall A \in \aF_t, B \in \aF^t, \quad \mathbb{P}(A \cap B \mid X_t) =  \mathbb{P}(A \mid X_t) \mathbb{P}(B \mid X_t).
	\]
\end{defn}

On fait ensuite l'identification commode : soit $X$ une variable aléatoire et $\aG \subset \aA$ deux tribus. On dit que
\[
	\E(X \mid \aG) = \E(X \mid \aA)
\]
si pour tout $A \in \aA$, on a
\[
	\E(\E(X \mid \aG)\mathbb{1}_A) = \E(\E(X \mid \aA)\mathbb{1}_A).
\]
Cela ne veut pas dire que $\E(X \mid \aA)$ est $\aG$-mesurable. Dans ce cas, on a la caractérisation de la propriété de Markov par la proposition suivante.

\begin{prop}{}{}
	Les proposition suivantes sont équivalentes :
	\begin{enumerate}
		\item $X$ est un processus de Markov
		\item Soit $Y$ une fonction $\aF^t$-mesurable bornée, alors $\E(Y \mid \aF_t) = \E(Y \mid X_t)$.
		\item $\forall t \le s$ et $f$ une application $\aE$-mesurable bornée, $\E(f\circ X_s \mid X_t) = \E(f\circ X_s \mid \aF_t)$.
	\end{enumerate}
\end{prop}

Pour construire, une chaîne de Markov dans un cadre très général, on peut utiliser une famille de probabilités de transition $(P_{t,s})$.

\begin{defn}{}{}
	Une famille de probabilités de transition sur $(E,\mathcal E)$ est une famille de noyaux
	$(P_{t,s})_{0\le t\le s}$
	telle que :

	\begin{enumerate}
		\item pour tous $0\le t\le s$, $P_{t,s}(x,\cdot)$ est une probabilité sur $(E,\mathcal E)$ ;

		\item pour tout $A\in\mathcal E$, l'application
		      \[
			      x\longmapsto P_{t,s}(x,A)
		      \]
		      est $\mathcal E$-mesurable ;

		\item la relation de Chapman--Kolmogorov est satisfaite :
		      \[
			      P_{t,s}(x,A)
			      =
			      \int_E P_{r,s}(y,A)\,P_{t,r}(x,dy),
			      \qquad 0\le t\le r\le s.
		      \]
	\end{enumerate}
\end{defn}

La relation de Chapman--Kolmogorov exprime la compatibilité temporelle des transitions. Elle permet de reconstruire les lois finies-dimensionnelles d'un processus de Markov. Plus précisément, sous des hypothèses usuelles de mesurabilité, une mesure initiale $\mu$ sur $(E,\mathcal E)$ et une famille de probabilités de transition $(P_{t,s})_{0\le t\le s}$ déterminent un processus de Markov dont les probabilités de transition sont données par les noyaux $P_{t,s}$.

\medskip

Dans le cas particulier où $P_{t,s}=P_{s-t}$,
les probabilités de transition ne dépendent que de l'écart temporel. Le processus est alors dit homogène en temps. La relation de Chapman--Kolmogorov devient
\[
	P_{t+s}=P_tP_s,
	\qquad t,s\ge0,
\]
c'est-à-dire que la famille $(P_t)_{t\ge0}$ forme un semi-groupe de Markov.
